\section*{الفصل \ref{ch:b3:photochemistry} --- الكيمياء الضوئية}

\begin{solution}{exo:b3:photochemistry:1}
$E = hc/\lambda = \qty{5.442e-19}{J}$؛ $0.100/\num{5.442e-19} = \num{1.84e17}$ فوتون في الثانية، أي
\qty{3.05e-7}{einstein.s^{-1}}.
\end{solution}

\begin{solution}{exo:b3:photochemistry:2}
الممتص: $\num{1.00e-7} \times 600 = \qty{6.00e-5}{einstein}$؛ $\Phi = \num{2.0e-5}/\num{6.0e-5} = 0.33$.
\end{solution}

\begin{solution}{exo:b3:photochemistry:3}
تتفاعل جزيئات أكثر بكثير من الفوتونات الممتصة: فكل جزيء \ce{Cl2} متحلل ضوئيًا يبدأ
سلسلة، $\ce{Cl + H2 -> HCl + H}$ و$\ce{H + Cl2 -> HCl + Cl}$، تدور آلاف الدورات قبل
الإنهاء.
\end{solution}

\begin{solution}{exo:b3:photochemistry:4}
$D_0(\ce{O2}) = 2 \times 246.79 = \qty{493.58}{kJ/mol}$؛ $\lambda = hcN_A/D_0 = 0.119627/493\,580 =
\qty{242.4}{nm}$. $\ce{O3 -> O2 + O}$: $246.79 - 145.35 = \qty{101.44}{kJ/mol}$، $\lambda = \qty{1179}{nm}$.
\end{solution}

\begin{solution}{exo:b3:photochemistry:5}
$[Z]/[E] = 22\,000 \times 0.11/(1500 \times 0.40) = 4.0$: أي 80\,\% من $Z$ و20\,\% من $E$.
\end{solution}

\begin{solution}{exo:b3:photochemistry:6}
الهكسان-2-ون: ينشطر ثنائي الجذر 1،4 إلى إينول الأسيتون (ومنه الأسيتون) 
والبروبين، أو ينغلق في 1،2-ثنائي ميثيل سيكلوبيوتان-1-ول. وفي البنتان-2-ون يكون الكربون $\gamma$
هو \ce{CH3} الطرفي: فانشطار الرابطة \babelsublr{$\alpha$--$\beta$} يترك $\ce{CH2=CH2}$.
\end{solution}

\begin{solution}{exo:b3:photochemistry:7}
$K = \exp[(E_{\mathrm T}(\mathrm D) - 250)/RT]$ مع $RT = \qty{2.478}{kJ/mol}$: يعطي 289 القيمة $\eu^{15.7} = \num{7e6}$؛
ويعطي 253 القيمة $\eu^{1.2} = 3.4$؛ ويعطي 236 القيمة $\eu^{-5.6} = \num{3.5e-3}$. فالأول وحده ينقل عند كل
لقاء تقريبًا.
\end{solution}

\begin{solution}{exo:b3:photochemistry:8}
$k_2 = \num{6.0e-34} \times (230/300)^{-2.6} = \qty{1.20e-33}{cm^6.s^{-1}}$؛ $k_2[\mathrm M][\ce{O2}] = \num{1.20e-33} \times
\num{4.0e17} \times \num{8.4e16} = \qty{40}{s^{-1}}$؛ $[\ce{O}]/[\ce{O3}] = \num{1.0e-3}/40 = \num{2.5e-5}$.
\end{solution}

\begin{solution}{exo:b3:photochemistry:9}
$k(\ce{Cl + O3}) = \num{2.8e-11}\eu^{-250/230} = \num{9.4e-12}$؛ $k(\ce{Cl + CH4}) = \num{6.6e-12}\eu^{-1240/230} = \num{3.0e-14}$
\unit{cm^3.s^{-1}}. والسرعتان من الرتبة الأولى الزائفة 54 و\qty{0.012}{s^{-1}}: والاحتمال $54/54.01 = 0.9998$.
\end{solution}

\begin{solution}{exo:b3:photochemistry:10}
$k = \num{2.07e-12}\eu^{-1400/298.15} = \qty{1.89e-14}{cm^3.s^{-1}}$؛ $[\ce{O3}] = \num{8.0e-3} \times 1.5/\num{1.89e-14} =
\qty{6.3e11}{cm^{-3}}$. ويحتوي الهواء عند \qty{1}{bar} و\qty{298}{K} على $p/k_BT = \qty{2.43e19}{cm^{-3}}$: أي \qty{26}{ppb}. وفي الليل
لا يعود \ce{NO2} يتحلل ضوئيًا، ويُتلف NO الأوزون حتى يُستهلك أحدهما.
\end{solution}

\begin{solution}{exo:b3:photochemistry:11}
$E = hc/\lambda = \qty{6.347e-19}{J}$؛ $0.050/\num{6.347e-19} = \num{7.88e16}$ فوتون في الثانية $= \qty{1.31e-7}{einstein.s^{-1}}$.
الممتص: $\times(1 - 10^{-0.5}) = \times0.684$، أي \qty{8.95e-8}{einstein.s^{-1}}؛ والسرعة $0.25 \times \num{8.95e-8} =
\qty{2.24e-8}{mol.s^{-1}}$، أي \qty{7.5e-6}{mol.L^{-1}.s^{-1}} في \qty{3.0}{mL}.
\end{solution}

\begin{solution}{exo:b3:photochemistry:12}
$I_0 = \num{3.0e-6}/(1.25 \times 60) = \qty{4.0e-8}{einstein.s^{-1}}$. فعند تحويل ضعيف يظل المقياس الأكتيني
يمتص كل الضوء، ولا ترشّحه النواتج الملوّنة.
\end{solution}

\begin{solution}{pb:b3:photochemistry:1}
\textbf{1.} $2 \times 246.79 = \qty{493.58}{kJ/mol}$.
\textbf{2.} $\lambda = 0.119627/493\,580 = \qty{242.4}{nm}$.
\textbf{3.} $246.79 - 145.35 = \qty{101.44}{kJ/mol}$؛ \qty{1179}{nm}.
\textbf{4.} امتصاص الأوزون في المرئي ضعيف؛ أما امتصاصه القوي في فوق البنفسجي
(نحو 200 إلى \qty{310}{nm}) فيزيل الإشعاع الذي يُتلف الحمض النووي والجلد.
\textbf{5.} الضوء الأقصر من \qty{242}{nm} يمتصه \ce{O2} نفسه في أثناء نزوله: فلا ينفذ منه إلا القليل
إلى ما دون الستراتوسفير العليا.
\textbf{6.} $\ce{O2 ->[h\nu] 2 O}$؛ $\mathrm{O + O_2 + M \to O_3 + M}$؛ $\ce{O3 ->[h\nu] O2 + O}$؛ $\ce{O + O3 -> 2 O2}$.
\textbf{7.} $k_2 = \qty{1.20e-33}{cm^6.s^{-1}}$؛ $k_2[\mathrm M] = \qty{4.79e-16}{cm^3.s^{-1}}$.
\textbf{8.} $k_4 = \num{8.0e-12}\eu^{-2060/230} = \qty{1.03e-15}{cm^3.s^{-1}}$.
\textbf{9.} $\num{1.0e-3}/(\num{4.79e-16} \times \num{8.4e16}) = \num{2.49e-5}$.
\textbf{10.} حصيلة الأكسجين الفردي $2j_1[\ce{O2}] = 2k_4[\ce{O}][\ce{O3}]$ مع $[\ce{O}] = j_3[\ce{O3}]/(k_2[\mathrm M][\ce{O2}])$ (انظر
المبرهنة).
\textbf{11.} $[\ce{O3}] = \num{8.4e16}\sqrt{\num{1.0e-11} \times \num{4.79e-16}/(\num{1.0e-3} \times \num{1.03e-15})} =
\qty{5.7e12}{cm^{-3}}$، أي نسبة خلط $\num{1.4e-5}$ (\qty{14}{ppm}).
\textbf{12.} $\num{2.49e-5} \times \num{5.7e12} = \qty{1.4e8}{cm^{-3}}$.
\textbf{13.} الدورات الحفزية (أكاسيد النيتروجين، وأكاسيد الهيدروجين، والكلور) تضيف قنوات
فقد للأكسجين الفردي.
\textbf{14.} $\ce{Cl + O3 -> ClO + O2}$، $\ce{ClO + O -> Cl + O2}$؛ والمحصلة $\ce{O + O3 -> 2 O2}$.
\textbf{15.} $1/(\num{9.44e-12} \times \num{5.72e12}) = \qty{0.019}{s}$.
\textbf{16.} $1/(\num{4.033e-11} \times \num{1.423e8}) = \qty{174}{s}$.
\textbf{17.} $174/0.0185 = \num{9.4e3}$.
\textbf{18.} الكلور كله تقريبًا في صورة ClO: الفقد $2 \times \num{5.74e-3} \times \num{1.0e7} = \qty{1.1e5}{cm^{-3}.s^{-1}}$،
مقابل $2k_4[\ce{O}][\ce{O3}] = \qty{1.7e6}{cm^{-3}.s^{-1}}$ لخطوة تشابمان: فالدورة تضيف
نحو 7\,\%.
\textbf{19.} $\ce{ClO + O}$، الخطوة البطيئة (\qty{174}{s}).
\textbf{20.} $53.8/(53.8 + 0.012) = 0.99978$.
\textbf{21.} $\num{5.74e-3}/(\num{5.74e-3} + \num{1.41e-13} \times \num{1.0e9}) = 0.976$.
\textbf{22.} $p = 0.976$.
\textbf{23.} $0.976/0.024 = 40$.
\textbf{24.} على السحب، يحرر التفاعل $\ce{HCl + ClONO2 -> Cl2 + HNO3}$ الكلور، ويبقى حمض النتريك
في الجسيمات: فدون \ce{NO2} يقترب $p_2$ من 1 وتصير السلسلة
طويلة جدًا؛ وتُتلف دورةٌ عبر ثنائي ClO الأوزون دون ذرات O.
\textbf{25.} \textbf{نحو 40 جزيء أوزون لكل ذرة كلور قبل الاحتجاز} (في هذا
النموذج لارتفاع \qty{30}{km})؛ وكل ذرة تتحرر من خزاناتها مرات كثيرة.
\end{solution}
